\section{Conditional column entropy and projection ranks}\label{tra:conditional}

We now retain information after the entire row array is known. Under a
law supported on compatible pairs, the column permutations can remain
correlated even after this conditioning. The first lemma bounds that
conditional correlation in terms of the separate conditional column
deficits. The estimate retains conditional information that the unconditional
marginal comparison in Section~\ref{tra:duality} does not record. Its projection consequence
also follows from the weighted compatibility theorem; the proof here
shows how it follows from the conditional estimate itself.

Throughout this section $T_n$ is the sweep from Section~\ref{sec:prelim}.
Traces and Schatten norms are unnormalized and include regular
multiplicities. The geometry is balanced, with both side lengths within
a factor two of $\sqrt n$.

\subsection{Conditional information and projection overlap}\label{tra:conditional:compatibility}
Write $n=rc$, $m=\sqrt n$, with $m/2\le r,c\le2m$, and let
$H=(S_c)^r$ and $K=(S_r)^c$ be the row and column groups.
A pair $(h,u)\in H\times K$ is compatible when the row move $h$
followed by the column move $u$ can be performed in the opposite order,
that is, $uh\in HK$. Equivalently, every original column has its
entries sent to distinct destination rows. Regarding cells after the
row move as edges between original and intermediate columns, the
column permutations then give a proper edge coloring by the $r$
destination rows.

\begin{lemma}[Entropy after the row array is revealed]
\label{lem:conditional-column-entropy}
Let $Q$ be any probability law on compatible pairs $(h,u)$, and write
$u_k$ for the permutation in intermediate column $k$. Let $H_Q$ denote
Shannon entropy in nats under $Q$. Define
\[
 F=\mathcal D(Q_h\|U_H),\qquad
 J=\sum_{k=1}^c\mathbb E_h\mathcal D(Q_{u_k\mid h}\|U_{S_r}),\qquad
 I=\sum_{k=1}^cH_Q(u_k\mid h)-H_Q(u\mid h).
\]
For all sufficiently large $n$,
\begin{equation}
 I\ge n-O(n^{.57})-O((F+J)/\log n).
 \label{tra:conditional:eq4}
\end{equation}
The constants are absolute, uniformly over $Q$.
\end{lemma}
To compare this with unconditional column marginals, put
$D_C=\sum_k\mathcal D(Q_{u_k}\|U_{S_r})$ and
$S=\sum_k I_Q(h;u_k)$, where $I_Q$ denotes mutual information. Then
\[
 J=D_C+S,\qquad
 \sum_kH_Q(u_k)-H_Q(u\mid h)=I+S.
\]
Thus the lemma retains the conditional information $I$ itself;
replacing it by the latter expression would include the columns'
separate information about the row array.
\begin{proof}\label{tra:conditional:conditional-entropy-proof}
Fix \(h\), and expose intermediate columns according to independent continuous uniform priorities between 0 and 1. Inside each column expose edge colors in a fixed order. For a current edge let \(p\) be the color law conditional only on the previously exposed colors \textbf{in that column} (and \(h\)). If \(R\) colors remain in the column, mark as heavy colors those with
\[
                               p_i>m^{.04}/R,
\]
and write \(w\) for the light mass. The relative entropy of \(p\) from uniform on the \(R\) remaining colors is at least
\[
                   (1-w)(.04\log m-1).
\]
Indeed the heavy terms contribute at least \((1-w).04\log m\), and by the log-sum inequality the rest contribute at least \(w\log w\). The sum over edges of the expected relative entropies here, also averaged over \(h\), is exactly \(J\).

Call an edge bad if the multiplicity \(b\) of its (original column, intermediate column) pair is greater than \(m^{.04}\). Averaging over \(h\), the expected number of bad edges is \(O((F+1)/\log n)\). To see this, under uniform row permutations the number \(T\) of bad edges has \(\mathbb E\exp(c_0 T\log m)\le2\) with some absolute \(c_0>0\) for sufficiently large \(m\). In detail, to count configurations with \(T=l>0\), list the column pairs of high multiplicity and the rows producing those occurrences. For \(a\le l/m^{.04}\) such pairs, listing them costs at most \((c^2)^a\), and their sizes can be given in at most \(2^l\) ways. Listing the sets of rows costs at most
\[
                            (er/m^{.04})^l
\]
by the binomial bound. The probability of the specified images is at most \((e/c)^l\), using \((c)_x=c(c-1)\cdots(c-x+1)\ge(c/e)^x\) in every row. This proves the exponential estimate by summing (since \(r,c\) differ only by a bounded factor). Now use the standard entropy variational inequality with \(F\).

By the entropy chain rule, each contribution to \(I\) in the exposure procedure is the information the earlier columns give about the current edge color beyond the local (within-column) history. We detail a lower bound on the average contribution for a nonbad edge. Fix the local history for this calculation. By also classifying the current color as light or heavy, a lower bound is \(w\) times that mutual information calculated under the condition of a light color (discard if \(w=0\)). This uses that the classification is a function of the color given the local history, and column order is independent. Under the light condition, the local color probabilities are the light part of \(p\), normalized. Once earlier columns are seen, colors used there at this original column are forbidden. Thus relative entropy from the local light distribution is at least minus log of its mass still allowed, by the support bound for relative entropy.

Average this last estimate for a column priority \(t\), with the other priorities random. Even given the entire valid coloring with current color light, any color at the original column that is used in a \emph{different} intermediate column is forbidden with probability \(t\), since the priorities are independent. The \(b\) colors used within the current column at the same original column have normalized local light mass at most \(b m^{.04}/(Rw)\). Jensen's inequality for the minus log therefore bounds our information contribution below, on average, by
\[
                w\int_0^1 -\log\left(1-t+t\,\min\{1,bm^{.04}/(Rw)\}\right)\,dt .
\]
The bound indeed concerns information with the column order known throughout (and is averaged over it). Since \(\int_0^1-\log(1-t)\,dt=1\), this is at least
\[
                         w-O\left(\frac{m^{.08}}R\log(em)\right)
\]
for a nonbad edge. For example this follows directly by integrating, using a loss \(O(x\log(e/x))\) inside the factor \(w\) when the minimum in braces is \(x\); if the minimum is 1 the stated looser bound also suffices. Within any column the \(1/R\) sum is \(O(\log(em))\). The total losses of this latter form are \(O(m^{1.08}\log^2(em))=O(n^{.57})\). The bounds on expected bad edges and summed expected heavy mass thus give \eqref{tra:conditional:eq4}.

\end{proof}

\begin{proposition}\label{tra:conditional:crossing}
\label{tra:conditional:crossing-statement}\label{e:conditional-rank-crossing}
For the balanced grid above, let $A,D$ be nonzero orthogonal
projections in the group algebras of $H,K$, respectively, each a tensor
product over its individual lines. Let $q,s$ be their ranks in the
regular representations of $H,K$. For all sufficiently large $n$,
in the regular representation of $S_n$,
\begin{equation}
 \|AD\|_4^4\le(qs)^{1+C/\log n}\exp(O(n^{.57})).
 \label{tra:conditional:eq2}
\end{equation}
\end{proposition}
\begin{proof}\label{tra:conditional:crossing-proof}
First sample $h,u$ independently with product laws across their lines
and densities, relative to uniform on $H,K$, bounded by $q,s$.
We show that their compatibility probability is at most
\begin{equation}
 \exp\{-n+O(n^{.57})+(C/\log n)\log(qs)\}.
 \label{tra:conditional:eq3}
\end{equation}
If this probability is $e^{-L}>0$, condition on compatibility and
use $F,J,I$ from Lemma~\ref{lem:conditional-column-entropy} for the
conditioned law. The chain rule against the original product law,
and then the density caps against uniform, give
\[
 L\ge I,\qquad L\ge F+J+I-\log(qs).
\]
Write the losses in \eqref{tra:conditional:eq4} as
$E+\delta(F+J)$, where $E=O(n^{.57})$ and
$\delta=O(1/\log n)$. Substitution gives
\[
 L\ge n-E-\delta(F+J),\qquad
 L\ge n-E+(1-\delta)(F+J)-\log(qs).
\]
Taking $(1-\delta)$ times the first bound and $\delta$ times the
second proves \eqref{tra:conditional:eq3}. A zero probability needs
no argument.

Apply Lemma~\ref{lem:coefficient-compatibility} to the individual line
projections forming $A,D$. Their squared coefficient masses equal
their line regular ranks. The normalized squared densities have the
same ranks as caps, so the products of masses and caps on the two
sides are $q,s$. Consequently
\[
 \|AD\|_4^4=\operatorname{Tr}(ADAD)
 \le\frac{n!}{|H||K|}\,qs\,
       \mathbb P(\text{compatibility}).
\]
The probability has precisely the independent product laws covered by
\eqref{tra:conditional:eq3}. Since
$\log(n!/(|H||K|))=n+O(m\log m)$, the leading $n$ cancels and
\eqref{tra:conditional:eq2} follows.
\end{proof}
\subsection{Harmonic cancellation at small levels}
\begin{lemma}\label{tra:conditional:sparse}
\label{tra:conditional:sparse-statement}
Let $V_\lambda$ be the irreducible representation indexed by
$\lambda\vdash n$, with first row $n-k$. For sufficiently large $n$,
the sweep on $V_\lambda$ satisfies
\begin{equation}
               \|T_n(\lambda)\|_{\rm op} \le n^{-.005 k},
                     \qquad 1\le k\le n^{.61}.                         \label{tra:conditional:eq5}
\end{equation}
\end{lemma}
\begin{proof}\label{tra:conditional:sparse-proof}
\label{tra:conditional:sparse-forest-proof}
By branching, $V_\lambda$ occurs on ordered $k$-slot injections but
not on $(k-1)$-slot injections. It therefore lies in the orthogonal
complement of functions of proper subtuples. We will bound the sweep
on that complement by cancellation of isolated paths.

For input and output \(k\)-tuples of distinct positions \(x,y\), the paths in one pass are uniquely prescribed (change the input bits in sequence to the output bits). Make a graph on paths, connecting any two touching the same switch anywhere. For \(V\subset[1,k]\), let \(M_V\) be \(n^{|V|}\) times the probability that the paths indexed by \(V\) are realized. This is zero for an incompatible collection and otherwise \(2^e\), where \(e\) counts switches used by two members, since coins on different switches are independent. In projecting the pass to the part orthogonal to functions of proper subtuples of the \(k\) cards, we can replace the kernel entries \(n^{-k}M_{[1,k]}\) by
\[
                      n^{-k}\sum_{V}(-1)^{k-|V|}M_V.
\]
Terms depending on a proper subtuple do not contribute, and the irreducible in question is in this orthogonal part. If there is any isolated path the sum vanishes because adding that path does not change \(M_V\).

Bound the squared Hilbert-Schmidt norm of this replacement. By Cauchy-Schwarz, up to factors \(\exp(O(k))\), it suffices to bound, uniformly for \(V\),
\[
          n^{-2k}\sum_{x,y\ \mathrm{as\ above}}
                        M_V^2 {\bf1}_{\{\text{no isolated path}\}}.
\]
One factor \(M_V\) allows the following path sampling upper bound on this expression:
\begin{itemize}
\item For \(V\), sample a uniform input injection and run one pass, keeping those paths.
\item For the rest, sample independent uniform input-output slot pairs.
\item Take expected \(2^e\) on the no-isolated event, with \(e\) counting shared switches among \(V\).
\end{itemize}
This is an upper bound because the normalization factor from requiring a distinct input for \(V\) is at most 1, and restrictions on the independent pairs may be dropped.

Condition first on the complete switch settings for that pass, so \(V\) samples without replacement from \(n\) routes using those settings. For a prescribed rooted spanning forest of the touch graph with \(a\) components, the probability of the required touches is bounded by
\[
                              (4d/n)^{k-a}.
\]
Indeed expose outward from roots. For a child among \(V\), given the parent path, at most \(2d\) of the full routes can touch the path, with at least \(n-k\) routes still to choose from. For an independent input-output pair the marginal at each layer is uniform, giving the same bound by a union bound. Moreover, within a component of size \(l\), shared switches among compatible pass routes (here the \(V\) routes) number at most \(l\log_2(l)/2\). This deterministic bound follows by splitting after the first layer into output-bit halves and inducting over remaining layers: for child group sizes \(l_1,l_2\) of a collection of routes the first layer contributes at most \(\min(l_1,l_2)\), at most half the binary entropy gain
\((l_1+l_2)\log_2(l_1+l_2)-\sum_i l_i\log_2 l_i\), by concavity.

Thus, taking a union bound with spanning forest witnesses for the actual components, all of sizes \(l\ge2\), the weight \(2^e\) costs a factor at most \(\prod_l l^{l/2}\) (product over components). Forests can be counted by selecting roots and specifying a parent for each nonroot, at most \(2^k k^{k-a}\) choices with a given \(a\). These bounds give at most \(n^{-.02 k}\), even including any of the \(\exp(O(k))\) factors above, for sufficiently large \(n\). To make the exponent check explicit, before those exponential factors the cost for a given number of components, allowing the worst sizes, is bounded using factors per component
\[
                          (4d k/n)^{l-1} l^{l/2}.
\]
For \(l\ge10\), the power saving from \((n/k)^{l-1}\) after paying \(l^{l/2}\) is at least \(n^{(.39\cdot .9-.305)l}\) since \(k\le n^{.61}\). For smaller sizes \(l^{l/2}\) costs only exponential constants. Polylog factors per vertex and summing over \(a\) do not affect the stated looser bound \(n^{-.02 k}\). Taking square roots gives \eqref{tra:conditional:eq5}.
\end{proof}
\subsection{A bounded moment exponent}
\begin{theorem}\label{tra:conditional:moment}
There is an absolute $p_*\ge4$ such that
\begin{equation}
              {\rm Tr}|T_n|^{p_*}\ \le\ 1+1/8                         \label{tra:conditional:eq1}
\end{equation}
for every power of two $n$. 
\end{theorem}
\begin{proof}[Proof of Theorem~\ref{tra:conditional:moment}]
\label{tra:conditional:bounded-moment-induction}
For the sweep recursion take $n=2^d$, $r=2^{\lceil d/2\rceil}$
and $c=2^{\lfloor d/2\rfloor}$. Fix $p_0'\ge10000$ and suppose the child sweeps satisfy
$\operatorname{Tr}|T_r|^{p_0'},\operatorname{Tr}|T_c|^{p_0'}\le1+1/8$.
We first obtain a rough parent bound, uniformly in $p_0'$.
Write $X=K_C$ for the column sweep and $Y=K_R$ for the row sweep,
so the chronological convention gives $T_n=XY$. We will show
\begin{equation}
             {\rm Tr}|XY|^p\le\exp(O(n^{.58})),
       \qquad p=p_0'(1+C_1/\log n),                                     \label{tra:conditional:eq6}
\end{equation}
for a sufficiently large absolute \(C_1\).

Lemma~\ref{lem:positive-power-comparison} gives
\[
            {\rm Tr}|XY|^p \le
                    \big\|\, |X|^{p/4}|Y^*|^{p/4}\,\big\|_4^4
                   \quad(p\ge4).
\]
We explain the summation bounding the right hand side to keep \eqref{tra:conditional:eq6} uniform in \(p\). Decompose each local positive factor (for a single row or column) by bins of singular values of the local pass, grouping singular values whose \(p_0'\)-powers lie in
\((e^{-j-1},e^{-j}]\), integer \(j\ge0\); zero parts can be omitted. Expand into products by choosing one bin per row or column. This works within the group algebras by spectral calculus. In the local regular representation the projection rank for such a bin is at most \(2 e^{j+1}\), by the child moment bound. For any term in the expansion of \(|X|^{p/4}|Y^*|^{p/4}\), let \(x\) be the sum of all its \(j\)'s. Its fourth norm power is bounded, by \eqref{tra:conditional:eq2} for the two support projections and the operator norm bounds on the factors, by
\[
 \exp\{-x p/p_0'+(1+C/\log n)(x+O(m))+O(n^{.57})\}.
\]
Taking \(C_1>C+1\), fourth roots can be summed by the triangle inequality at additional cost at most \((O(\log n))^{O(m)}\), by geometric series. This proves \eqref{tra:conditional:eq6}.

The projection calculation has supplied only the rough bound
\eqref{tra:conditional:eq6}. To reduce it to a fixed small remainder,
we separate the sparse levels just bounded from representations whose
regular multiplicity is exponentially large.

Partitions neither having first row of length at least \(n-n^{.61}\), nor first column that long, have dimension at least
\[
                                \exp(n^{.60})
\]
for large \(n\). Here is one direct check. If the longest row and column are shorter than \(2n^{.61}\), the hook formula suffices immediately, all hooks being at most \(4n^{.61}\). Otherwise (transpose if necessary) use a subdiagram with first row of length about \(2n^{.61}\) and about \(n^{.61}\) boxes below (round down, discarding excess). By branching it suffices to lower bound dimension there. After filling the initial portion of the first row of length equal to the number of lower boxes, the remaining first row entries and lower boxes can be interleaved freely in a standard tableau (with any one standard ordering of lower boxes), giving the claimed bound.

The first-row levels $1\le k\le n^{.61}$ have at most $2^k$
shapes at each $k$, each of dimension at most $n^k$. Therefore
\eqref{tra:conditional:eq5} bounds their regular contribution at
any exponent $u\ge10000$ by
\[
 \sum_{1\le k\le n^{.61}}2^k n^{2k-.005uk}=o(1).
\]
The shapes with first column at least $n-n^{.61}$ are annihilated
by Lemma~\ref{lem:annihilation}, since this height exceeds $n/2$
for large $n$. Together these form the first class in
Lemma~\ref{lem:balanced-moment-closure}; its contribution is at most
$1/16$ above an absolute threshold, uniformly in $u$.

For the remaining class take $H_n=n^{.60}$ from the dimension bound,
$E_n=C_2n^{.58}$ from \eqref{tra:conditional:eq6}, and
$\delta_d=1/\log n$. Then
\[
 (1+\delta_d)E_n-\delta_dH_n
 =(1+1/\log n)C_2n^{.58}-n^{.60}/\log n\longrightarrow-\infty.
\]
The common lemma applies with singular convention $\nu=1$,
$\varepsilon=1/8$, $P=10000$, and
$\alpha_d=1+C_1/\log n$. The combined multiplier is
$1+O(1/d)$. It supplies one finite $p_*$ and proves
\eqref{tra:conditional:eq1}, including the finite initial sizes.
\label{tra:conditional:fourier-completion}
The invariant in \eqref{tra:conditional:eq1} has singular exponent
$p_*$ and remainder $1/8$. Its conversion to forward sweep powers is
therefore the case $2b\ge p_*$ of
Proposition~\ref{e:nonnormal-moment-conversion}.
\end{proof}
